\section{Математические модели в биологии и биоинформатике}
\label{sec:models-in-biology-and-bioinformatics}

Математическое моделирование в биологии используется для описания биологических процессов с помощью математических формул. Это позволяет изучать динамику популяций, распространение заболеваний, взаимодействие клеток и организмов, работу генов и белков, а также делать прогнозы.

В биоинформатике математические модели применяются, например, для анализа последовательностей ДНК и РНК, поиска сходства между генами, восстановления филогенетических связей и обработки биологических данных.

Основные задачи можно условно разделить на несколько типов.

\subsection{Моделирование роста популяции}

Пусть $N(t)$ --- численность популяции в момент времени $t$. Простая модель неограниченного роста имеет вид

\[
\frac{dN}{dt}=rN,
\]

где $r$ --- коэффициент роста. Она описывает ситуацию, когда ресурсов достаточно и численность популяции растёт экспоненциально.

Если учитывать ограниченность ресурсов, используется логистическая модель:

\[
\frac{dN}{dt}
=
rN\left(1-\frac{N}{K}\right),
\]

где $K$ --- ёмкость среды. В этом случае рост постепенно замедляется, и численность популяции стремится к значению $K$.

\subsection{Моделирование взаимодействия видов}

Например, отношения хищник--жертва можно описать моделью Лотки--Вольтерры:

\[
\begin{cases}
	\dfrac{dx}{dt}=ax-bxy,\\[6pt]
	\dfrac{dy}{dt}=-cy+dxy.
\end{cases}
\]

Здесь $x$ --- численность жертв, а $y$ --- численность хищников.

Первый член $ax$ описывает размножение жертв, а слагаемые, содержащие произведение $xy$, описывают взаимодействие хищника и жертвы. Модель позволяет исследовать колебания численности двух видов.

\subsection{Распространение инфекционного заболевания}

Популяцию можно разделить на восприимчивых к болезни $S$, инфицированных $I$ и выздоровевших $R$.

В простой SIR-модели процесс описывается системой:

\[
\begin{cases}
	\dfrac{dS}{dt}=-\beta SI,\\[6pt]
	\dfrac{dI}{dt}=\beta SI-\gamma I,\\[6pt]
	\dfrac{dR}{dt}=\gamma I.
\end{cases}
\]

Здесь $\beta$ характеризует скорость заражения, а $\gamma$ --- скорость выздоровления.

Такая модель позволяет оценивать динамику эпидемии, например максимальное число одновременно заболевших и скорость распространения инфекции.

\subsection{Задачи биоинформатики: сравнение биологических последовательностей}

Для двух последовательностей ДНК или белков

\[
A=a_1a_2\dots a_n,
\qquad
B=b_1b_2\dots b_m
\]

можно искать оптимальное выравнивание.

Каждому совпадению, несовпадению и пропуску назначается определённый балл. Тогда задача заключается в нахождении такого выравнивания, которое максимизирует суммарную оценку:

\[
\max S(A,B).
\]

Подобные задачи используются для определения сходства генов или белков, а также для поиска возможного общего происхождения биологических последовательностей.


\newpage